\section{Preguntas frecuentes}

\begin{longtable}{@{}p{0.38\textwidth}p{0.52\textwidth}@{}}
\toprule
\textbf{Pregunta} & \textbf{Respuesta}\\
\midrule
Olvidé mi contraseña & Pídale a Administración TI que la restablezca; el
  portal no envía claves por correo.\\
La cuenta está bloqueada & Se bloquea 15 minutos tras cinco intentos
  fallidos. Administración TI puede desbloquearla antes.\\
No veo una pantalla & Su rol no la incluye. El menú muestra solo lo
  autorizado.\\
Cambié un dato y el PDF no cambia & El PDF se genera desde la copia
  congelada. Hay que corregir la ficha, lo que crea una versión nueva.\\
El estudiante no aparece al buscar & Revise el año escolar activo y pruebe
  con el RUN sin puntos o con el IPE.\\
Falta un documento & El expediente no bloquea la matrícula: se marca
  pendiente y se completa después.\\
¿Puedo borrar una matrícula? & No se borra: se anula con motivo y se conserva
  el historial.\\
¿Quién puede ver la salud del estudiante? & Secretaría, TENS, UTP, Dirección
  e Inspectoría. Docente y Convivencia no.\\
\bottomrule
\end{longtable}

\section{Glosario}

\begin{longtable}{@{}p{0.26\textwidth}p{0.64\textwidth}@{}}
\toprule
\textbf{Término} & \textbf{Significado}\\
\midrule
Folio & Número único que identifica una ficha o un certificado.\\
Versión & Cada corrección de una ficha genera una versión nueva.\\
Copia congelada & Datos guardados al cerrar la ficha; de ahí sale el PDF.\\
Autorización & Permiso firmado por el apoderado (imagen, convivencia,
  paracetamol).\\
Prescripción & Orden médica que respalda la administración de un medicamento.\\
Expediente & Conjunto de documentos del estudiante y su estado.\\
Validar & Revisión de Dirección que deja la matrícula revisada y registrada.\\
Anular & Dejar sin efecto una matrícula conservando su historial.\\
RUN / IPE & Cédula nacional / identificador provisorio de estudiantes
  extranjeros.\\
Rol & Perfil de trabajo que define qué pantallas y acciones tiene una persona.\\
\bottomrule
\end{longtable}

\section{Soporte}

\begin{itemize}[leftmargin=*,itemsep=0.2em]
\item \textbf{Cuentas, claves y permisos}: Administración TI del colegio.
\item \textbf{Datos de un estudiante}: Secretaría.
\item \textbf{Matrículas y certificados}: Dirección.
\item \textbf{Fallas del sistema}: quien administra el servidor del portal.
\end{itemize}

Al pedir ayuda, indique la pantalla, lo que intentaba hacer, el mensaje que
apareció y la hora. Si es un error de datos, agregue el RUN del estudiante.
